\documentclass[10pt,a4paper]{amsart}
\usepackage[margin=19mm]{geometry}
\usepackage{amsmath,amssymb,mathtools}
\usepackage[scr=rsfs,cal=euler]{mathalfa}
\usepackage{xfrac}
\usepackage[unicode,hidelinks,bookmarks=false]{hyperref}
\newcommand{\ie}{i.e.\ }
\newcommand{\cf}{cf.\ }
\newcommand{\resp}{respectively\ }
\newcommand{\Subsection}{\subsection}
\providecommand{\nobreakdash}{\nobreak\hbox{-}}
\setlength{\emergencystretch}{2em}
\multlinegap0pt
\usepackage{amssymb}
\usepackage{enumerate}
\usepackage{tikz}
\usepackage{mathtools}
\newtheorem{lemma}{Lemma}
\title{Asymptotic v-number of graded families of ideals and the Newton-Okounkov region}
\author{Mousumi Mandal, Partha Phukan}
\date{Source: arXiv:2603.08838v3 (CC BY 4.0)}
\begin{document}
\maketitle

\noindent\emph{Source and licence.} Asymptotic v-number of graded families of ideals and the Newton-Okounkov region. Version arXiv:2603.08838v3 (CC BY 4.0), distributed by arXiv under the Creative Commons Attribution 4.0 International licence, http://creativecommons.org/licenses/by/4.0/, as recorded in the arXiv licence metadata for this exact version and shown on the abstract page https://arxiv.org/abs/2603.08838v3. Full licence notice: \texttt{licenses/arxiv-2603.08838-CC-BY-4.0.txt}.\par\smallskip

\noindent\textit{Editorial scope.} Counted unit: the article's lemma ArtinReeslemma (source0.tex lines 1330--1335). The article's complete original proof of it is delivered with it (source0.tex lines 1337--1360, one \texttt{proof} environment of 24 lines). No mathematics is written, repaired or re-derived: the statement and the proof are the article's own lines, in the article's own notation, and no retained line is substituted or re-flowed.  No further source-local context is needed: every cross-reference of every retained line resolves inside the two delivered spans.  Nothing was omitted from any retained span, so the package carries no omission record.  No retained line contains a citation, so the wrapper carries no bibliography and no bibliography entry.  No retained line contains a figure environment, an image-inclusion command or drawing code, so no image is delivered and the package declares no asset dependency.  Editorial formatting only, in addition: an AMS \texttt{amsart} preamble that restates the article's own declarations and macros used by the retained lines, namely the environments \texttt{lemma} on separate counters, together with the standard packages those lines need.  The retained environments print in this document as Lemma 1 in order of appearance, whereas the article prints the same items with the numbering of its own full document.  The complete native source of the article as distributed in the arXiv e-print for this exact version is bundled unchanged as \texttt{sources/196032/source0.tex}.
\begin{lemma}\label{ArtinReeslemma}
Let $R$ be a Noetherian ring, $I\subseteq R$ an ideal, and let $a\in R$ be a non-zero divisor. Then there exists an integer $s\ge 1$ such that, for every $k\ge 0$,
\[
(I^{k+s}:_R a)\subseteq I^k.
\]
\end{lemma}

\begin{proof}
By the Artin-Rees Lemma, there exists an integer $s\ge 1$ such that
\(
I^n\cap aR = I^{n-s}(I^s\cap aR)
\)
for all $n\ge s$. Set $n=k+s$. Then
\(
I^{k+s}\cap aR = I^k(I^s\cap aR)
\)
for all $k\ge 0$.

\noindent For every integer $m\ge 0$, we have
\(
I^m\cap aR = a(I^m:_R a).
\)
Now,
\[
a(I^{k+s}:_R a) = I^{k+s}\cap aR = I^k(I^s\cap aR) = I^k\cdot a(I^s:_R a) \subseteq aI^k.
\]
Since $a$ is a non-zero divisor, it follows that
\(
(I^{k+s}:_R a)\subseteq I^k.
\)
\end{proof}

\end{document}
